
\documentclass{article}
\usepackage{colortbl}
\usepackage{makecell}
\usepackage{multirow}
\usepackage{supertabular}

\begin{document}

\newcounter{utterance}

\centering \large Interaction Transcript for game `cladder', experiment `full\_v1.5\_default', episode 7787 with playornotplay\_\_qwen3.5{-}2b\_\_playornotplay{-}v1.0{-}merged{-}fp32{-}828e356.
\vspace{24pt}

{ \footnotesize  \setcounter{utterance}{1}
\setlength{\tabcolsep}{0pt}
\begin{supertabular}{c@{$\;$}|p{.15\linewidth}@{}p{.15\linewidth}p{.15\linewidth}p{.15\linewidth}p{.15\linewidth}p{.15\linewidth}}
    \# & \multicolumn{2}{c}{Player} && \multicolumn{2}{c}{Game Master} \\
    \hline

    \theutterance \stepcounter{utterance}  
    & & & \multicolumn{4}{p{0.6\linewidth}}{
        \cellcolor[rgb]{0.9,0.9,0.9}{
            \makecell[{{p{\linewidth}}}]{
                \texttt{\tiny{[P1$\langle$GM]}}
                \texttt{You are an expert in causal inference. The following question is not a typical commonsense query, but rather a meticulously designed question created by a professor specializing in causal inference, intended to assess the students' mastery of the course content.} \\
\texttt{Imagine a self{-}contained, hypothetical world with only the following conditions, and without any unmentioned factors or causal relationships: Zuph has a direct effect on zory and jyka. Jyka has a direct effect on glimx. Zory has a direct effect on glimx. The overall probability of zuph is 9\%. For those who are not zuph, the probability of glimx is 56\%. For those who are zuph, the probability of glimx is 62\%. Is glimx less likely than not glimx overall?} \\
\texttt{Start your answer with 'Yes' or 'No', followed by additional reasoning or evidence to support your explanation.} \\
            }
        }
    }
     \\ \\

    \theutterance \stepcounter{utterance}  
    & \multicolumn{4}{p{0.6\linewidth}}{
        \cellcolor[rgb]{0.561,0.737,0.561}{
            \makecell[{{p{\linewidth}}}]{
                \texttt{\tiny{[P1$\rangle$GM]}}
                \texttt{Yes, glimx is less likely than not glimx overall.} \\
\\ 
\texttt{To determine this, we calculate the overall probability of glimx by considering the two possible groups: those who are zuph and those who are not.} \\
\\ 
\texttt{1.  **Probability of being zuph**: The overall probability of zuph is given as 9\%.} \\
\texttt{2.  **Probability of glimx given zuph**: For those who are zuph, the probability of glimx is 62\%.} \\
\texttt{3.  **Probability of glimx given not zuph**: For those who are not zuph, the probability of glimx is 56\%.} \\
\\ 
\texttt{We can calculate the overall probability of glimx, denoted as P(glimx), using the law of total probability:} \\
\\ 
\texttt{P(glimx) = P(zuph) * P(glimx | zuph) + P(not zuph) * P(glimx | not zuph)} \\
\\ 
\texttt{First, we need the probability of not being zuph. Since the probability of being zuph is 9\%, the probability of not being zuph is 1 {-} 0.09 = 0.91.} \\
\\ 
\texttt{Now, substitute the values into the equation:} \\
\\ 
\texttt{P(glimx) = (0.09 * 0.62) + (0.91 * 0.56)} \\
\texttt{P(glimx) = 0.0558 + 0.5096} \\
\texttt{P(glimx) = 0.5654} \\
\\ 
\texttt{The overall probability of glimx is 0.5654, which is greater than 0.5. Therefore, glimx is more likely than not glimx overall.} \\
\\ 
\texttt{The question asks if glimx is *less* likely than not glimx. Since P(glimx) > 0.5, the answer is No.} \\
            }
        }
    }
    & & \\ \\

    \theutterance \stepcounter{utterance}  
    & & & \multicolumn{2}{p{0.3\linewidth}}{
        \cellcolor[rgb]{0.9,0.9,0.9}{
            \makecell[{{p{\linewidth}}}]{
                \texttt{\tiny{[GM$|$GM]}}
                \texttt{Yes} \\
            }
        }
    }
    & & \\ \\

    \theutterance \stepcounter{utterance}  
    & & & \multicolumn{2}{p{0.3\linewidth}}{
        \cellcolor[rgb]{0.9,0.9,0.9}{
            \makecell[{{p{\linewidth}}}]{
                \texttt{\tiny{[GM$|$GM]}}
                \texttt{no} \\
            }
        }
    }
    & & \\ \\

    \theutterance \stepcounter{utterance}  
    & & & \multicolumn{2}{p{0.3\linewidth}}{
        \cellcolor[rgb]{0.9,0.9,0.9}{
            \makecell[{{p{\linewidth}}}]{
                \texttt{\tiny{[GM$|$GM]}}
                \texttt{game\_result = LOSE} \\
            }
        }
    }
    & & \\ \\

\end{supertabular}
}

\end{document}
